\section{Results}
\label{sec:results}

\subsection{Controlled Reimplementation}

\begin{table}[t]
\centering
\caption{BTCV controlled reimplementation. Published values are shown only for
calibration; our paired analyses use the local runs. MAC denotes
multiply--accumulate operations under our profiler.}
\label{tab:baseline}
\resizebox{\columnwidth}{!}{
\begin{tabular}{lrrrrr}
\toprule
Model & Dice $\uparrow$ & HD95 $\downarrow$ & MAC & Params & Lat.\\
 & & & (G) & (M) & (ms)\\
\midrule
Published \fullmodel{} & .7974 & -- & 631.97 & 53.00 & --\\
Published \effimodel{} & .7925 & 10.12 & 51.47 & 3.16 & --\\
\midrule
\fullmodel{}, seed 0 & .7918 & 9.04 & 578.74 & 53.007 & 40.6\\
\effimodel{}, seed 0 & .7700 & 14.41 & 41.06 & 2.955 & 7.5\\
\effimodel{}, seed 1 & .7755 & 10.17 & 41.06 & 2.955 & 7.6\\
\bottomrule
\end{tabular}}
\end{table}

Table~\ref{tab:baseline} shows that the full model closely approaches the
published Dice (a 0.56-point difference), whereas the two efficient runs remain
1.70--2.25 points below the published efficient model.  The two-seed efficient
mean is 0.7728, and seed 1 nearly matches the published HD95.  The local
\fullmodel{}--\effimodel{} Dice gap (1.63--2.18 points) is larger than the
published 0.49-point gap; we do not attribute its full magnitude to decoder
compression.  Efficiency trends reproduce strongly: \effimodel{} uses
$14.1\times$ fewer MACs, $17.9\times$ fewer parameters, $5.3$--$5.4\times$ lower
latency, and roughly $7.7\times$ lower peak inference memory than \fullmodel{}.

\subsection{Errors and Uncertainty Are Anatomically Concentrated}

\begin{figure}[t]
\centering
\includegraphics[width=0.96\columnwidth]{figures/O1_organ_error.png}\\[-1mm]
\textbf{(a) Per-organ error.}\\[1mm]
\includegraphics[width=0.96\columnwidth]{figures/O6_entropy_evolution.png}\\[-1mm]
\textbf{(b) Entropy over training.}
\caption{Anatomical and temporal structure of difficulty (\effimodel{}, seed 1).
Error concentrates at boundaries and small organs; entropy contracts during
training but a residual difficult subset persists.}
\label{fig:anatomy}
\end{figure}

The premise for selective allocation is that difficulty is not spatially
uniform.  Mean boundary error is $0.1892$ versus $0.0481$ in the eroded
interior, a ratio of $3.93$ (Fig.~\ref{fig:anatomy}a).  Small and thin
structures are disproportionately hard: the organ-size/difficulty Spearman
correlation is $\rho=-0.54$, and entropy retains strong association with organ
difficulty after controlling for log volume (partial $r=0.81$), so size alone
does not explain the signal.

\begin{figure}[t]
\centering
\includegraphics[width=0.86\columnwidth]{figures/O2_entropy.png}
\caption{Voxel entropy distribution (\effimodel{}, seed 1; log-count).
Uncertainty is spatially sparse: the median voxel entropy is near zero and only
$1.18\%$ of voxels exceed $0.5$.}
\label{fig:entropy}
\end{figure}

Uncertainty is likewise sparse: the median voxel entropy is near zero and only
$1.18\%$ of voxels exceed $0.5$ (Fig.~\ref{fig:entropy}).  Difficulty also
evolves systematically---mean entropy falls from $0.0279$ at 5k iterations to
$0.0166$, $0.0126$, $0.0123$, and $0.0104$ at 10k, 20k, 30k, and 45k
(Fig.~\ref{fig:anatomy}b)---but training removes global ambiguity without
eliminating the concentrated residual.  Most voxels are confidently easy; any
benefit from a stronger decoder must come from the sparse remainder.

\subsection{The Full Decoder's Benefit Is Small and Bidirectional}

\begin{table}[t]
\centering
\caption{Transition-count observations across efficient-decoder seeds. Rates are
fractions of evaluated voxels. The top-20\% recovery is a voxel oracle
(Sec.~\ref{sec:study}), not the deployable region result.}
\label{tab:transitions}
\begin{tabular}{lrr}
\toprule
Metric & seed 0 & seed 1\\
\midrule
Positive transition rate & .00339 & .00290\\
Negative transition rate & .00209 & .00225\\
Net transition rate & .00130 & .00065\\
Subject entropy--net $r$ & .665 & .646\\
Bootstrap 95\% CI & [.171,.996] & [.158,.994]\\
Top-20\% positive recovery & .864 & .864\\
\bottomrule
\end{tabular}
\end{table}

The full decoder is not uniformly better.  For seed 0 it improves $0.339\%$ of
voxel predictions while degrading $0.209\%$; for seed 1 the values are $0.290\%$
and $0.225\%$ (Table~\ref{tab:transitions}).  The net label-transition benefit is
therefore only $0.065$--$0.130\%$ of voxels, largely offset by negative
transitions.  Although this label-level statistic is not identical to Dice, it
shows why positive and negative transitions must be reported together: reporting
positive transitions alone would overstate the decoder's benefit roughly
twofold.

\subsection{The Benefit Is Spatially Predictable}

\begin{figure}[t]
\centering
\includegraphics[width=0.96\columnwidth]{figures/O5_decoder_gain.png}\\[-1mm]
\textbf{(a) Net transition vs.\ entropy.}\\[1mm]
\includegraphics[width=0.96\columnwidth]{figures/O9_opportunity_corrected.png}\\[-1mm]
\textbf{(b) Region-level opportunity curve.}
\caption{Predictability (seed 1). Net transitions rise with entropy (a); ranking
contiguous $16^3$ blocks by entropy recovers net transition utility far above a
matched-random selector (b).}
\label{fig:opportunity}
\end{figure}

The subject-level entropy--net-transition correlation is positive for both
seeds, with a subject-bootstrap lower bound above zero
(Table~\ref{tab:transitions}, Fig.~\ref{fig:opportunity}a).  We convert this into
a deployable test with the contiguous-region selector of
Section~\ref{sec:study}.  At the predeclared $20\%$ block budget---which executes
$29.1\%$ of the foreground-union volume after halo expansion---entropy-ranked
regions attain a net utility of $0.225$ versus $0.043$ for random selection, a
roughly fivefold gain, with a paired $95\%$ lower bound of $0.056>0$
(Fig.~\ref{fig:opportunity}b).  The advantage is significant at every budget of
$10\%$ or more and vanishes only at $5\%$.  Crucially, the selected regions
capture all positive transitions but also $77.5\%$ of the negative ones, so the
honest headline is the \emph{net} gain, not positive recovery.  A voxel-level
oracle that ranks individual voxels and counts only positive transitions places
$86.4\%$ of them in the top $20\%$ (Table~\ref{tab:transitions}); we report it
solely as a non-deployable upper bound.

Finally, we compare cheap routing signals against held-out net transitions.
Confidence and entropy behave almost identically ($r=0.663$ and $0.655$);
Monte-Carlo dropout is not a valid comparator for this checkpoint because the
architecture has no active dropout, so its stochastic variance is effectively
zero.  We therefore retain confidence as a required baseline and do not claim
entropy is uniquely optimal.

\todo{Add the matched Swin UNETR replication and the predeclared FeTA/MSD cells;
current results are 3D UX-Net on BTCV only.}
